\documentclass[11pt, oneside]{article} 
\usepackage{geometry}
\geometry{letterpaper} 
\usepackage{graphicx}
	
\usepackage{amssymb}
\usepackage{amsmath}
\usepackage{parskip}
\usepackage{color}
\usepackage{hyperref}

\graphicspath{{/Users/telliott/Github-Math/figures/}}
% \begin{center} \includegraphics [scale=0.4] {gauss3.png} \end{center}

\title{Cycloid}
\date{}

\begin{document}
\maketitle
\Large

%[my-super-duper-separator]

Imagine a bicycle with one tire marked at a particular point on the rim, say with fluorescent paint or a small light.  Time starts at $t = 0$ with that point $P$ in contact with the $x$ axis at $(0,0)$.  Then the bike rolls to our right.  As the tire rotates the fixed point $P$ on the rim traces a curve
\begin{center} \includegraphics [scale=0.6] {cycloid.png} \end{center}

We want to find equations that give the point $P$ as a function of time.  We will parametrize the curve, yielding parametric equations $x(t)$, $y(t)$.

The second diagram  shows the angle through which the wheel has turned as $\theta$, but we will use $t$ for $\theta$ here.  

\begin{center} \includegraphics [scale=0.5] {cycloid2.png} \end{center}
The $x$ displacement of the vertical straight down from the center of the tire is just $at$, where $a$ is the radius of the wheel, it is equal to the arc on the circumference of the wheel from the point which is currently in contact with the ground, going around up to $P$.

It is reasonably easy to derive the desired parametric equations, using vectors, especially once you know the answer.  For $x$, we have the vector that goes from $(0,0)$ to the contact point with the ground.  As indicated in the figure, that is $at$.  

We need to subtract the distance $a \ sin\ t$ from that.  Basically the rationale is that the motion is a standard parametric circle which has been rotated by $90$ degrees clockwise and then inverted.  The rotation changes cosine to sine, and the inversion brings the subtraction.

\begin{center} \includegraphics [scale=0.35] {cycloid3.png} \end{center}

It's easier to see for $t < \pi/2$, but it is true always.  Check some other values of $t$ like $\pi$ or $3\pi/2$ to confirm.  This is the usual circular motion.

For $y$, we have a constant factor of $a$ above the $x$ axis, then the additional displacement is $-a \ cos \ t$.  So for $t=0$ we have the additional displacement is -a (we were on the ground), for $t=\pi/2$ it is zero, and for $t=\pi$ it is plus $a$ for a total of $2a$.

The parametric equations are then
\[ x(t) = at - a \sin t \]
\[ y(t) = a - a \cos t \]

Once we know some calculus, we can get various expressions like the slope of the curve

Taking derivatives:
\[ x'(t) = a - a \cos t \]
\[ y'(t) = a \sin t  \]
We can get $dy/dx$, what we usually call $y'$, by simple division:
\[ y' = \frac{dy}{dx} = \frac{a \sin t }{a - a \cos t} \]
\[ = \frac{\sin t}{1 - \cos t} \]
Here, we will do what we can without calculus.

Simmons uses half-angle formula to make something simpler, like so:
\[ \sin t = 2 \sin \frac{t}{2} \cos \frac{t}{2} \]
\[ \cos t = \cos^2 \frac{t}{2} - \sin^2 \frac{t}{2} \]
\[ = 1 - 2 \sin^2 \frac{t}{2} \]
So the ratio is
\[ y' = \frac{ 2 \sin \frac{t}{2} \cos \frac{t}{2}}{2 \sin^2 \frac{t}{2}} = \cot \frac{t}{2} \]

\subsection*{Roberval}
The cycloid curve was not known to the ancients, as far as our sources have anything to say.  The first mention is around 1500, and by the time of Galileo a century later, it was something he thought about quite a bit and encouraged others to study.  

Galileo is said to have used Archimedes method of cutting out shapes and weighing them to show that the area was approximately $3$ times that of the generating circle.  But he was apparently not prepared for it to be \emph{exactly} equal, so was skeptical.

The cycloid became very popular in the decades before Newton.  Roberval came up with a clever application of Cavlieri's principle to find the area under the curve.

The great idea is to draw a second curve, called the \emph{companion curve}, by sliding values derived from a half-circle to add them to the $x$-values on the cycloid curve, as shown.

The areas marked by horizontal lines are equal, by Cavalieri's principle.
But that total extra area is just $\pi a^2 / 2$.

\url{https://maa.org/sites/default/files/pdf/cmj_ftp/CMJ/January%202010/3%20Articles/3%20Martin/08-170.pdf}

\begin{center} \includegraphics [scale=0.2] {cycloid_comp.png} \end{center}

Notice that the curve $AQD$ apparently divides the rectangular area $ACDB$ in half, by symmetry.  Since the rectangle has width $\pi a$ and height $2a$, its area is $2 \pi a^2$, so one-half that is $\pi a^2$.

The area under the cycloid curve is then three-quarters of the total, which (between $0 \rightarrow \pi$) is $(3/2) \pi a^2$, and the area under one complete lobe is twice that or $3\pi a^2$.

Let's see if we can figure out a parametric equation for the companion curve.  $x(t)$ for the cycloid was $x(t) = at - a \sin t$.  The companion curve gets an additional length in the $x$-direction of $a \sin t$.  So $x(t) = at$.  That was easy!

If we want $y = f(x)$ we must do
\[ y = a - a \cos t \]
and then substitute $t = x/a$ so
\[ y = a(1 - \cos x/a) \]

which goes like the sine squared of a scaled version of $x$.

\begin{center} \includegraphics [scale=0.2] {cycloid5.png} \end{center}

\subsection*{slope}
Descartes gave the slope of the tangent line to the cycloid by the following construction:
\begin{center} \includegraphics [scale=0.2] {cycloid_slope.png} \end{center}
Given $P$ on the cycloid, draw $PE$ horizontally to find $E$ on the circle, then draw $PQ$ parallel to $EC$.  The tangent is perpendicular to $PQ$.

I like this construction a lot because it reminds me of what Roberval did with the parabola.  For any $t$, the distance moved horizontally is equal to the distance along the rim of the generating circle.  The tangent should add to or subtract from both values at the same rate.  

In other words, I expect that the tangent should be the average of the tangent to the circle at $E$ and the horizontal, namely one-half of the former.

If $C$ is taken to be $(0,0)$, the coordinates of $E$ are $(x',y)$, where
\[ x' = -a \sin t  \]
and the slope of $PQ \parallel EC$ will be 
\[ \frac{\Delta y}{\Delta x} = \frac{0 - (a - a \cos t)}{0 - (-a \sin t)} \]
\[ = \frac{\cos t - 1}{\sin t} = \cot t - \csc t \]
The slope of the tangent is the negative inverse
\[ \frac{\sin t}{1 - \cos t} \]
which matches.  Recall that we had before that $y' =  \cot t/2$
\begin{center} \includegraphics [scale=0.2] {cycloid_slope.png} \end{center}
which seems like a problem at first but $t \ne \angle ECQ$.  

If we had drawn the center of the circle as $O$, then $t$ would be $\angle EOC$, and since $\triangle EOC$ is isosceles
\[ t + 2 \cdot \angle OCE = 180 \]
\[ t/2 + \angle OCE = 90 \]
but $\angle OCE$ ($\angle DCE$) and $\angle ECQ$ are also complementary so
\[ t/2 = \angle ECQ = \angle EPQ \]

The angle made by the tangent ($\angle HPE$) is complementary, so its tangent is the cotangent of $t/2$.

Note also that $\angle EDC = \angle ECQ = \angle EPQ$, so it is also complementary to the tangent angle.

\end{document}